% !TEX program = lualatex
% Math 110 — chapter-by-chapter notes scaffold
% ============================================================
%  Math 110 — Chapter-by-Chapter Reading Notes
%  Axler, \emph{Linear Algebra Done Right} (4th ed.)
%  Keshav Ramamurthy — Fall 2026
% ------------------------------------------------------------
%  HOW TO USE (one growing document, one banner per chapter)
%
%    Find the chapter banner and the section heading below it,
%    then type your notes directly underneath:
%
%        Big idea: every Cauchy sequence converges iff ...
%
%    Structure the important ideas with the note environments
%    (numbered within the current chapter):
%
%        \begin{definition} ... \end{definition}
%        \begin{theorem} ... \end{theorem}
%        \begin{example}  ... \end{example}
%        \begin{remark}   ... \end{remark}
%        \begin{idea}     ... \end{idea}
%
%    Add a missing section:  \notesec{§1.4}{Title}
%    Add a chapter:          \chapterbanner{1}{Chapter 1}{Title}{description}
% ============================================================
\documentclass[11pt]{article}

\usepackage{fontspec}
\usepackage{microtype}
\usepackage{mathtools,amssymb,amsfonts,amsthm,bm}
\usepackage{enumitem}
\usepackage{booktabs,array,xcolor}
\usepackage{geometry,fancyhdr,setspace}
\usepackage{etoolbox}
\usepackage[most]{tcolorbox}
\usepackage[hidelinks]{hyperref}
\usepackage[nameinlink,noabbrev]{cleveref}

% ---------- Page and paragraph rhythm (matches lecture/solutions templates) ----------
\geometry{margin=1in,headheight=15pt}
\setstretch{1.12}
\setlength{\parindent}{0pt}
\setlength{\parskip}{0.55em plus 0.12em minus 0.08em}
\setlength{\abovedisplayskip}{0.8em plus 0.2em minus 0.1em}
\setlength{\belowdisplayskip}{0.8em plus 0.2em minus 0.1em}
\allowdisplaybreaks[3]
\setlist{leftmargin=*,topsep=0.35em,itemsep=0.15em,parsep=0pt}

% ---------- Palette (matches reference cheatsheet) ----------
\definecolor{ink}{HTML}{1E293B}
\definecolor{navy}{HTML}{183B56}
\definecolor{paperblue}{HTML}{E8F1F8}
\definecolor{gold}{HTML}{FFF3CD}
\definecolor{mint}{HTML}{E8F5E9}
\definecolor{muted}{HTML}{52606D}
\definecolor{rule}{HTML}{C9D5DF}
\color{ink}

% ---------- Metadata ----------
\newcommand{\studentname}{Keshav Ramamurthy}
\newcommand{\coursename}{Math 110}
\newcommand{\courselongname}{Abstract Linear Algebra}
\newcommand{\textbookname}{Axler, \emph{Linear Algebra Done Right} (4th ed.)}
\newcommand{\term}{Fall 2026}

% ---------- Core notation (same as ~/.config/nvim/textemplate.tex) ----------
\newcommand{\N}{\mathbb{N}}
\newcommand{\Z}{\mathbb{Z}}
\newcommand{\Q}{\mathbb{Q}}
\newcommand{\R}{\mathbb{R}}
\newcommand{\C}{\mathbb{C}}
\newcommand{\F}{\mathbb{F}}
\newcommand{\K}{\mathbb{K}}
\newcommand{\E}{\mathbb{E}}
\newcommand{\Prob}{\mathbb{P}}
\newcommand{\1}{\mathbf{1}}
\newcommand{\eps}{\varepsilon}
\newcommand{\dd}{\,\mathrm{d}}
\newcommand{\abs}[1]{\left\lvert#1\right\rvert}
\newcommand{\norm}[1]{\left\lVert#1\right\rVert}
\newcommand{\ip}[2]{\left\langle#1,#2\right\rangle}
\newcommand{\set}[1]{\left\{#1\right\}}
\newcommand{\ceil}[1]{\left\lceil#1\right\rceil}
\newcommand{\floor}[1]{\left\lfloor#1\right\rfloor}
\newcommand{\given}{\,\middle\vert\,}
\newcommand{\indep}{\perp\!\!\!\perp}
\newcommand{\wh}[1]{\widehat{#1}}
\newcommand{\deriv}[2]{\frac{\mathrm{d}#1}{\mathrm{d}#2}}
\newcommand{\pderiv}[2]{\frac{\partial#1}{\partial#2}}
\newcommand{\lto}{\longrightarrow}
\newcommand{\st}{\text{ such that }}
\DeclareMathOperator{\Aut}{Aut}
\DeclareMathOperator{\End}{End}
\DeclareMathOperator{\Gal}{Gal}
\DeclareMathOperator{\Hom}{Hom}
\DeclareMathOperator{\Ker}{ker}
\DeclareMathOperator{\im}{im}
\DeclareMathOperator{\rank}{rank}
\DeclareMathOperator{\nullity}{nullity}
\DeclareMathOperator{\Span}{span}
\DeclareMathOperator{\tr}{tr}
\DeclareMathOperator{\diag}{diag}
\DeclareMathOperator{\spec}{spec}
\DeclareMathOperator{\ord}{ord}
\DeclareMathOperator{\supp}{supp}
\DeclareMathOperator{\diam}{diam}
\DeclareMathOperator{\Var}{Var}
\DeclareMathOperator{\Cov}{Cov}
\DeclareMathOperator*{\argmin}{arg\,min}
\DeclareMathOperator*{\argmax}{arg\,max}

% ---------- Note environments (numbered by chapter) ----------
\newcounter{chap}
\newtheoremstyle{notestyle}{0.7em}{0.7em}{\itshape}{}{}{\bfseries}{.5em}{}
\theoremstyle{notestyle}
\newtheorem{theorem}{Theorem}[chap]
\newtheorem{lemma}[theorem]{Lemma}
\newtheorem{proposition}[theorem]{Proposition}
\newtheorem{corollary}[theorem]{Corollary}
\theoremstyle{definition}
\newtheorem{definition}[theorem]{Definition}
\newtheorem{example}[theorem]{Example}
\newtheorem{remark}[theorem]{Remark}

% ---------- Note machinery ----------
% Chapter banner: \chapterbanner{1}{Chapter 1}{Groups and Subgroups}{Fraleigh: ...}
\newtcolorbox{chapterbox}[2]{%
  enhanced,
  colback=paperblue, colframe=navy, boxrule=0.7pt, arc=2.5pt,
  left=10pt, right=10pt, top=12pt, bottom=8pt,
  boxed title style={colback=navy, colframe=navy, arc=2pt, boxrule=0pt,
    left=8pt, right=8pt, top=3pt, bottom=3pt},
  attach boxed title to top left={yshift=-2.4mm, xshift=5mm},
  title={\sffamily\bfseries\color{white}\large #1\quad #2}}
\newcommand{\chapterbanner}[4]{%
  \clearpage
  \setcounter{chap}{#1}%
  \setcounter{theorem}{0}%
  \phantomsection
  \addcontentsline{toc}{section}{#2\quad #3}%
  \begin{chapterbox}{#2}{#3}
    {\small\sffamily\color{muted}#4\par}
  \end{chapterbox}
  \medskip\nopagebreak}

% Section heading: \notesec{§1.4}{Title}  →  navy heading + thin rule
\newcommand{\notesec}[2]{%
  \par\goodbreak\medskip\nopagebreak
  {\sffamily\bfseries\color{navy}#1\quad #2}%
  \par\nobreak\vspace{1.5pt}
  {\color{rule}\hrule height 0.6pt}\par\nobreak\smallskip\nopagebreak
  \phantomsection
  \addcontentsline{toc}{subsection}{#1\quad #2}}

% Quick structure inside notes
\newenvironment{claim}{\par\smallskip\noindent\textit{Claim.}\ }{\par\smallskip}
\newenvironment{idea}{\par\smallskip\noindent\textit{Idea.}\ }{\par\smallskip}
\newenvironment{proofsketch}{\par\smallskip\noindent\textit{Proof sketch.}\ }{\par\smallskip}

% ---------- Header ----------
\pagestyle{fancy}
\fancyhf{}
\fancyhead[L]{\coursename}
\fancyhead[C]{\courselongname\ — chapter notes}
\fancyhead[R]{\studentname}
\fancyfoot[C]{\thepage}

\begin{document}

% ---------- Title ----------
\begin{center}
  {\Large\sffamily\bfseries\color{navy}\coursename\ -- \courselongname}\\[3pt]
  {\sffamily\color{muted} Chapter-by-chapter notes on the important ideas}\\[3pt]
  {\small\sffamily\color{muted}\textbookname\quad\textbullet\quad \term\quad\textbullet\quad \studentname}
\end{center}
\medskip
{\color{rule}\hrule height 0.8pt}
\bigskip

{\small
\begin{tcolorbox}[colback=paperblue, colframe=rule, boxrule=0.4pt, arc=2pt,
  left=7pt, right=7pt, top=5pt, bottom=5pt,
  title={\sffamily\bfseries\color{navy} Taking notes},
  colbacktitle=paperblue]
\sffamily Find the chapter banner, then type under the section heading as you read:
\begin{center}
\texttt{\textbackslash notesec\{1A\}\{Span and Linear Independence\}}
\end{center}
Record big ideas with \texttt{idea}, definitions with \texttt{definition}, theorems with
\texttt{theorem}, \texttt{lemma}, \texttt{proposition}, examples with \texttt{example}, and
warnings with \texttt{remark} --- all numbered within the chapter.
Add a missing section by copying a \texttt{\textbackslash notesec} slot;
add a chapter with \texttt{\textbackslash chapterbanner}.
\end{tcolorbox}
}

\medskip
\tableofcontents

\chapterbanner{1}{Chapter 1}{Vector Spaces}{Axler:\ \(\F^n\), vector space axioms, subspaces, sums and direct sums}

\notesec{1A}{\(\R^n\) and \(\C^n\)}
% Big idea:
% Key definitions \& theorems:
% Examples / tricks:
% Questions:

\notesec{1B}{Definition of Vector Space}
% Big idea:
% Key definitions \& theorems:
% Examples / tricks:
% Questions:

\notesec{1C}{Subspaces}
% Big idea:
% Key definitions \& theorems:
% Examples / tricks:
% Questions:

\chapterbanner{2}{Chapter 2}{Finite-Dimensional Vector Spaces}{Axler:\ span, linear independence, bases, dimension}

\notesec{2A}{Span and Linear Independence}
% Big idea:
% Key definitions \& theorems:
% Examples / tricks:
% Questions:

\notesec{2B}{Bases}
% Big idea:
% Key definitions \& theorems:
% Examples / tricks:
% Questions:

\notesec{2C}{Dimension}
% Big idea:
% Key definitions \& theorems:
% Examples / tricks:
% Questions:

\chapterbanner{3}{Chapter 3}{Linear Maps}{Axler:\ linear maps, null space and range, matrices, invertibility and isomorphisms, products and quotients, duality}

\notesec{3A}{Vector Space of Linear Maps}
% Big idea:
% Key definitions \& theorems:
% Examples / tricks:
% Questions:

\notesec{3B}{Null Spaces and Ranges}
% Big idea:
% Key definitions \& theorems:
% Examples / tricks:
% Questions:

\notesec{3C}{Matrices}
% Big idea:
% Key definitions \& theorems:
% Examples / tricks:
% Questions:

\notesec{3D}{Invertibility and Isomorphisms}
% Big idea:
% Key definitions \& theorems:
% Examples / tricks:
% Questions:

\notesec{3E}{Products and Quotients of Vector Spaces}
% Big idea:
% Key definitions \& theorems:
% Examples / tricks:
% Questions:

\notesec{3F}{Duality}
% Big idea:
% Key definitions \& theorems:
% Examples / tricks:
% Questions:

\chapterbanner{4}{Chapter 4}{Polynomials}{Axler:\ zeros, the division algorithm, factorization over \(\C\) and over \(\R\)}

\notesec{4.1}{Zeros of Polynomials}
% Big idea:
% Key definitions \& theorems:
% Examples / tricks:
% Questions:

\notesec{4.2}{Division Algorithm for Polynomials}
% Big idea:
% Key definitions \& theorems:
% Examples / tricks:
% Questions:

\notesec{4.3}{Factorization of Polynomials over \(\C\)}
% Big idea:
% Key definitions \& theorems:
% Examples / tricks:
% Questions:

\notesec{4.4}{Factorization of Polynomials over \(\R\)}
% Big idea:
% Key definitions \& theorems:
% Examples / tricks:
% Questions:

\chapterbanner{5}{Chapter 5}{Eigenvalues and Eigenvectors}{Axler:\ invariant subspaces, eigenvalues, the minimal polynomial, upper-triangular matrices, diagonalizability, commuting operators}

\notesec{5A}{Invariant Subspaces}
% Big idea:
% Key definitions \& theorems:
% Examples / tricks:
% Questions:

\notesec{5B}{The Minimal Polynomial}
% Big idea:
% Key definitions \& theorems:
% Examples / tricks:
% Questions:

\notesec{5C}{Upper-Triangular Matrices}
% Big idea:
% Key definitions \& theorems:
% Examples / tricks:
% Questions:

\notesec{5D}{Diagonalizable Operators}
% Big idea:
% Key definitions \& theorems:
% Examples / tricks:
% Questions:

\notesec{5E}{Commuting Operators}
% Big idea:
% Key definitions \& theorems:
% Examples / tricks:
% Questions:

\chapterbanner{6}{Chapter 6}{Inner Product Spaces}{Axler:\ inner products and norms, orthonormal bases and Gram--Schmidt, orthogonal complements, minimization}

\notesec{6A}{Inner Products and Norms}
% Big idea:
% Key definitions \& theorems:
% Examples / tricks:
% Questions:

\notesec{6B}{Orthonormal Bases}
% Big idea:
% Key definitions \& theorems:
% Examples / tricks:
% Questions:

\notesec{6C}{Orthogonal Complements and Minimization Problems}
% Big idea:
% Key definitions \& theorems:
% Examples / tricks:
% Questions:

\chapterbanner{7}{Chapter 7}{Operators on Inner Product Spaces}{Axler:\ adjoints, self-adjoint and normal operators, the spectral theorems, positive operators, isometries, SVD}

\notesec{7A}{Self-Adjoint and Normal Operators}
% Big idea:
% Key definitions \& theorems:
% Examples / tricks:
% Questions:

\notesec{7B}{The Spectral Theorem}
% Big idea:
% Key definitions \& theorems:
% Examples / tricks:
% Questions:

\notesec{7C}{Positive Operators}
% Big idea:
% Key definitions \& theorems:
% Examples / tricks:
% Questions:

\notesec{7D}{Isometries, Unitary Operators, and Matrix Factorization}
% Big idea:
% Key definitions \& theorems:
% Examples / tricks:
% Questions:

\notesec{7E}{Singular Value Decomposition}
% Big idea:
% Key definitions \& theorems:
% Examples / tricks:
% Questions:

\notesec{7F}{Consequences of Singular Value Decomposition}
% Big idea:
% Key definitions \& theorems:
% Examples / tricks:
% Questions:

\chapterbanner{8}{Chapter 8}{Operators on Complex Vector Spaces}{Axler:\ generalized eigenvectors, nilpotent operators, generalized eigenspace decomposition, Jordan form, trace}

\notesec{8A}{Generalized Eigenvectors and Nilpotent Operators}
% Big idea:
% Key definitions \& theorems:
% Examples / tricks:
% Questions:

\notesec{8B}{Generalized Eigenspace Decomposition}
% Big idea:
% Key definitions \& theorems:
% Examples / tricks:
% Questions:

\notesec{8C}{Consequences of Generalized Eigenspace Decomposition}
% Big idea:
% Key definitions \& theorems:
% Examples / tricks:
% Questions:

\notesec{8D}{Trace:\ A Connection Between Matrices and Operators}
% Big idea:
% Key definitions \& theorems:
% Examples / tricks:
% Questions:

\chapterbanner{9}{Chapter 9}{Multilinear Algebra and Determinants}{Axler:\ bilinear and quadratic forms, alternating multilinear forms, determinants, tensor products}

\notesec{9A}{Bilinear Forms and Quadratic Forms}
% Big idea:
% Key definitions \& theorems:
% Examples / tricks:
% Questions:

\notesec{9B}{Alternating Multilinear Forms}
% Big idea:
% Key definitions \& theorems:
% Examples / tricks:
% Questions:

\notesec{9C}{Determinants}
% Big idea:
% Key definitions \& theorems:
% Examples / tricks:
% Questions:

\notesec{9D}{Tensor Products}
% Big idea:
% Key definitions \& theorems:
% Examples / tricks:
% Questions:

\end{document}
